% ============================================================
% methodology.tex — Sección V: Metodología Propuesta
% Estado: Conceptual (sin implementación detallada)
% ============================================================

\section{Metodología}
\label{sec:methodology}

El desarrollo de \sirca{} sigue la metodología \textit{Design Science Research} (DSR)~\cite{hevner2004design}, que proporciona un marco riguroso para la creación y evaluación de artefactos tecnológicos innovadores. La investigación se estructura en las siguientes fases:

\begin{enumerate}[label=\textbf{Fase \arabic*:}]
  \item \textbf{Fundamentación teórica y análisis de requisitos} (fase actual): Estudio del estado del arte, formulación del problema, definición de los fundamentos teóricos y diseño conceptual de la arquitectura.

  \item \textbf{Construcción del corpus y prototipado:} Adquisición de literatura biomédica sobre plantas medicinales peruanas mediante APIs científicas, procesamiento documental, vectorización e indexación.

  \item \textbf{Implementación del prototipo:} Desarrollo de los cinco componentes de \sirca{} utilizando el stack tecnológico propuesto, con integración de los módulos de retrieval, generación y memoria.

  \item \textbf{Evaluación experimental:} Evaluación cuantitativa mediante las métricas definidas en la Sección~\ref{sec:experiments}, comparación con baselines y análisis de resultados.
\end{enumerate}

El corpus biomédico se construirá mediante adquisición programática desde PubMed, CrossRef, Semantic Scholar y Europe PMC, utilizando consultas centradas en las especies medicinales peruanas seleccionadas (\textit{Uncaria tomentosa}, \textit{Lepidium meyenii}, \textit{Croton lechleri}, \textit{Minthostachys mollis}, \textit{Erythroxylum coca}) y sus propiedades farmacológicas. Los documentos adquiridos se someterán a un pipeline de preprocesamiento que incluye extracción de texto, segmentación semántica (Ecuación~\ref{eq:chunking}), extracción de entidades biomédicas mediante \scispacy{} y vectorización con embeddings biomédicos.
